% GENERATED FILE. Edit canonical lessons, then run npm run generate:books.
% canonical-source-hash: fnv1a64:f7a3e7646c561895
% canonical-lessons: BN-C382-nishtha, BN-C382-odhyobosay, BN-C382-puron-kora, BN-C382-bojay-rakha, BN-C382-songshodhon-kora

\chapter{Keeping It Up}
\label{ch:bn-keeping-it-up}
\hlchaptermodality{382}

\begin{chapteropening}
\textbf{By the end of this chapter:} \emph{I can say what it takes to keep going and finish what I start.}

Complete the last lesson of chapter 382: I can say what it takes to keep going and finish what I start.
\end{chapteropening}

\section[\texorpdfstring{\bn{নিষ্ঠা} (nishṭhā)}{নিষ্ঠা (nishṭhā)}]{\bn{নিষ্ঠা} (nishṭhā) — dedication; devotion (to work)}

\label{lesson:BN-C382-nishtha}



\begin{quote}
Before the new one: say the Bengali for skill, then the Bengali for to practise.
\end{quote}

\subsection*{You'll want to know: \bn{নিষ্ঠা}}
\textbf{\bn{নিষ্ঠা}} — \emph{nishṭhā} — \textquotedblleft{}dedication; devotion (to work)\textquotedblright{}.

\begin{quote}
\bn{নিষ্ঠা} \bn{থাকলে} \bn{তুমি} \bn{কঠিন} \bn{কাজও} \bn{শিখতে} \bn{পারবে।}
\end{quote}

\begin{quote}
\emph{nishṭhā thākle tumi kôṭhin kājo shikhte pārbe.}
\end{quote}

\begin{quote}
With dedication, you can learn even a hard skill.
\end{quote}

\begin{cousinweb}
\textbf{\bn{নিষ্ঠা}} (\emph{nishṭhā}) is Sanskrit \emph{niṣṭhā}, firmness and steady devotion (Monier-Williams): staying with a task out of care for it.
\end{cousinweb}

\begin{grammarlens}[title={What you've built}]
One more word for saying what happens if.
\end{grammarlens}

\subsection*{Your turn}
\begin{itemize}
\raggedright
  \item \emph{Say it:} \emph{nishṭhā}
  \item \emph{Say it:} \emph{nishṭhā}, once more
  \item \emph{Say it:} \emph{nishṭhā thākle tumi kôṭhin kājo shikhte pārbe.} — With dedication, you can learn even a hard skill.
  \item \emph{Recall:} say the Bengali for skill, then the Bengali for to practise, then say \emph{nishṭhā} again
\end{itemize}

\begin{tcolorbox}[breakable,colback=teal!4,colframe=teal!35!black,title={Before you move on}]
What does \bn{নিষ্ঠা} mean? (\textquotedblleft{}Dedication; devotion (to work)\textquotedblright{}.) Say it once more.
\end{tcolorbox}
\section[\texorpdfstring{\bn{অধ্যবসায়} (ôdhyôbôsāy)}{অধ্যবসায় (ôdhyôbôsāy)}]{\bn{অধ্যবসায়} (ôdhyôbôsāy) — perseverance}

\label{lesson:BN-C382-odhyobosay}



\begin{quote}
Before the new one: say the Bengali for to practise, then the Bengali for dedication.
\end{quote}

\subsection*{You'll want to know: \bn{অধ্যবসায়}}
\textbf{\bn{অধ্যবসায়}} — \emph{ôdhyôbôsāy} — \textquotedblleft{}perseverance\textquotedblright{}.

\begin{quote}
\bn{অধ্যবসায়} \bn{থাকলে} \bn{সাফল্য} \bn{আসবেই।}
\end{quote}

\begin{quote}
\emph{ôdhyôbôsāy thākle sāphôlyô āsbei.}
\end{quote}

\begin{quote}
If you persevere, success will surely come.
\end{quote}

\begin{cousinweb}
\textbf{\bn{অধ্যবসায়}} (\emph{ôdhyôbôsāy}) is Sanskrit \emph{adhyavasāya}, resolve or effort of mind (Monier-Williams). Bengali uses it for the patience to keep trying.
\end{cousinweb}

\begin{grammarlens}[title={What you've built}]
One more word for saying what happens if.
\end{grammarlens}

\subsection*{Your turn}
\begin{itemize}
\raggedright
  \item \emph{Say it:} \emph{ôdhyôbôsāy}
  \item \emph{Say it:} \emph{ôdhyôbôsāy}, once more
  \item \emph{Say it:} \emph{ôdhyôbôsāy thākle sāphôlyô āsbei.} — If you persevere, success will surely come.
  \item \emph{Recall:} say the Bengali for to practise, then the Bengali for dedication, then say \emph{ôdhyôbôsāy} again
\end{itemize}

\begin{tcolorbox}[breakable,colback=teal!4,colframe=teal!35!black,title={Before you move on}]
What does \bn{অধ্যবসায়} mean? (\textquotedblleft{}Perseverance\textquotedblright{}.) Say it once more.
\end{tcolorbox}
\section[\texorpdfstring{\bn{পূরণ} \bn{করা} (pūrôn kôrā)}{পূরণ করা (pūrôn kôrā)}]{\bn{পূরণ} \bn{করা} (pūrôn kôrā) — to fulfil; to fill in (a form)}

\label{lesson:BN-C382-puron-kora}



\begin{quote}
Before the new one: say the Bengali for dedication, then the Bengali for perseverance.
\end{quote}

\subsection*{You'll want to know: \bn{পূরণ} \bn{করা}}
\textbf{\bn{পূরণ} \bn{করা}} — \emph{pūrôn kôrā} — \textquotedblleft{}to fulfil; to fill in (a form)\textquotedblright{}.

\begin{quote}
\bn{এই} \bn{ফর্মটা} \bn{পূরণ} \bn{করলে} \bn{তুমি} \bn{ছাড়} \bn{পাবে।}
\end{quote}

\begin{quote}
\emph{ei phôrmṭā pūrôn kôrle tumi chhāṛ pābe.}
\end{quote}

\begin{quote}
If you fill in this form, you'll get a discount.
\end{quote}

Past (he/she …): \bn{তিনি} \bn{তাঁর} \bn{স্বপ্ন} \bn{পূরণ} \bn{করলেন}, \emph{tini tā̃r svôpnô pūrôn kôrlen}.

\begin{cousinweb}
\textbf{\bn{পূরণ} \bn{করা}} (\emph{pūrôn kôrā}) is Sanskrit \emph{pūraṇa}, filling or completing (Monier-Williams), with \bn{করা}: to fill a form, a need or a dream.
\end{cousinweb}

\begin{grammarlens}[title={What you've built}]
One more word for saying what happens if.
\end{grammarlens}

\subsection*{Your turn}
\begin{itemize}
\raggedright
  \item \emph{Say it:} \emph{pūrôn kôrā}
  \item \emph{Say it:} \emph{pūrôn kôrā}, once more
  \item \emph{Say it:} \emph{ei phôrmṭā pūrôn kôrle tumi chhāṛ pābe.} — If you fill in this form, you'll get a discount.
  \item \emph{Recall:} say the Bengali for dedication, then the Bengali for perseverance, then say \emph{pūrôn kôrā} again
\end{itemize}

\begin{tcolorbox}[breakable,colback=teal!4,colframe=teal!35!black,title={Before you move on}]
What does \bn{পূরণ} \bn{করা} mean? (\textquotedblleft{}To fulfil; to fill in (a form)\textquotedblright{}.) Say it once more.
\end{tcolorbox}
\section[\texorpdfstring{\bn{বজায়} \bn{রাখা} (bôjāy rākhā)}{বজায় রাখা (bôjāy rākhā)}]{\bn{বজায়} \bn{রাখা} (bôjāy rākhā) — to maintain; to keep up}

\label{lesson:BN-C382-bojay-rakha}



\begin{quote}
Before the new one: say the Bengali for perseverance, then the Bengali for to fulfil.
\end{quote}

\subsection*{You'll want to know: \bn{বজায়} \bn{রাখা}}
\textbf{\bn{বজায়} \bn{রাখা}} — \emph{bôjāy rākhā} — \textquotedblleft{}to maintain; to keep up\textquotedblright{}.

\begin{quote}
\bn{এই} \bn{গতি} \bn{বজায়} \bn{রাখলে} \bn{আমরা} \bn{সময়মতো} \bn{পৌঁছাব।}
\end{quote}

\begin{quote}
\emph{ei gôti bôjāy rākhle āmrā sômôymôto poũchhābô.}
\end{quote}

\begin{quote}
If we keep up this pace, we'll arrive on time.
\end{quote}

Past (he/she …): \bn{তিনি} \bn{শান্তি} \bn{বজায়} \bn{রাখলেন}, \emph{tini shānti bôjāy rākhlen}.

\begin{cousinweb}
\textbf{\bn{বজায়} \bn{রাখা}} (\emph{bôjāy rākhā}) is usually traced to Persian \emph{ba-jā}, in its place, with \bn{রাখা} 'to keep': to keep a thing where it is, to maintain it.
\end{cousinweb}

\begin{grammarlens}[title={What you've built}]
One more word for saying what happens if.
\end{grammarlens}

\subsection*{Your turn}
\begin{itemize}
\raggedright
  \item \emph{Say it:} \emph{bôjāy rākhā}
  \item \emph{Say it:} \emph{bôjāy rākhā}, once more
  \item \emph{Say it:} \emph{ei gôti bôjāy rākhle āmrā sômôymôto poũchhābô.} — If we keep up this pace, we'll arrive on time.
  \item \emph{Recall:} say the Bengali for perseverance, then the Bengali for to fulfil, then say \emph{bôjāy rākhā} again
\end{itemize}

\begin{tcolorbox}[breakable,colback=teal!4,colframe=teal!35!black,title={Before you move on}]
What does \bn{বজায়} \bn{রাখা} mean? (\textquotedblleft{}To maintain; to keep up\textquotedblright{}.) Say it once more.
\end{tcolorbox}
\section[\texorpdfstring{\bn{সংশোধন} \bn{করা} (sôngshodhôn kôrā)}{সংশোধন করা (sôngshodhôn kôrā)}]{\bn{সংশোধন} \bn{করা} (sôngshodhôn kôrā) — to correct; to amend}

\label{lesson:BN-C382-songshodhon-kora}



\begin{quote}
Before the new one: say the Bengali for to fulfil, then the Bengali for to keep up.
\end{quote}

\subsection*{You'll want to know: \bn{সংশোধন} \bn{করা}}
\textbf{\bn{সংশোধন} \bn{করা}} — \emph{sôngshodhôn kôrā} — \textquotedblleft{}to correct; to amend\textquotedblright{}.

\begin{quote}
\bn{ভুল} \bn{থাকলে} \bn{প্রবন্ধটা} \bn{সংশোধন} \bn{করো।}
\end{quote}

\begin{quote}
\emph{bhul thākle prôbôndhôṭā sôngshodhôn kôro.}
\end{quote}

\begin{quote}
If there are mistakes, correct the essay.
\end{quote}

Past (he/she …): \bn{তিনি} \bn{আইনটা} \bn{সংশোধন} \bn{করলেন}, \emph{tini āinṭā sôngshodhôn kôrlen}.

\begin{cousinweb}
\textbf{\bn{সংশোধন} \bn{করা}} (\emph{sôngshodhôn kôrā}) is Sanskrit \emph{saṃśodhana}, thorough cleansing (Monier-Williams), with \bn{করা}. Bengali cleans out errors: to correct a text or amend a law.
\end{cousinweb}

\begin{grammarlens}[title={What you've built}]
One more word for saying what happens if.
\end{grammarlens}

\subsection*{Your turn}
\begin{itemize}
\raggedright
  \item \emph{Say it:} \emph{sôngshodhôn kôrā}
  \item \emph{Say it:} \emph{sôngshodhôn kôrā}, once more
  \item \emph{Say it:} \emph{bhul thākle prôbôndhôṭā sôngshodhôn kôro.} — If there are mistakes, correct the essay.
  \item \emph{Recall:} say the Bengali for to fulfil, then the Bengali for to keep up, then say \emph{sôngshodhôn kôrā} again
\end{itemize}

\begin{tcolorbox}[breakable,colback=teal!4,colframe=teal!35!black,title={Before you move on}]
What does \bn{সংশোধন} \bn{করা} mean? (\textquotedblleft{}To correct; to amend\textquotedblright{}.) Say it once more.
\end{tcolorbox}
